\section{Finite endpoints, infinity, and powers of the inverse}
\label{sec:endpoints}

The branch theorem, coefficient formula and parametric convergence proof only need $p\ne0$ and $u\to0^+$: the substitution $t=u^p$ and the marker theorem are applied at positive points regardless of the endpoint approached by $t$. Statements that use positivity explicitly, such as the exponent-cutoff corollary to the Rouch\'e majorant, retain their stated restrictions. This section makes the endpoint bookkeeping explicit.

\subsection{Local variables}

Let $f$ be given near $x_0$, and let $u\to0^+$ be the local variable of $x$:
\begin{equation}
 u=\begin{cases}x-x_0 & x\to x_0^+,\\ x_0-x & x\to x_0^-,\\ 1/x & x\to+\infty,\\ -1/x & x\to-\infty .\end{cases}
 \label{eq:local-u}
\end{equation}
Substituting into $f$ and expanding (Section~\ref{sec:forward}) gives a germ $y_0+au^p(1+\sum u^{\delta_i}B_i(\log u))$ in the sense of \eqref{eq:model}, where $y_0=\lim f$ is finite when $p>0$. For $p<0$ the target endpoint is $\pm\infty$ with the sign of $a$, while the additive offset $y_0$ in \eqref{eq:model} is set to $0$; any constant term is a correction with gap $-p$. The local variable of $y$ is
\begin{equation}
 w=\begin{cases}\ \ \,y-y_0 & a>0,\ p>0,\\ -(y-y_0) & a<0,\ p>0,\\ \ \ \,1/|y| & p<0 .\end{cases}
 \label{eq:local-w}
\end{equation}
The uniformizer $z=(v/a)^{1/p}$ tends to $0^+$ in all cases. Its relation to $w$ is $w=|a|z^p$ for $p>0$ and $w=z^{|p|}/|a|$ for $p<0$.

\subsection{The theorem at an arbitrary endpoint}

\begin{theorem}[Signed leading power]
\label{thm:endpoints}
Theorems~\ref{thm:branch}, \ref{thm:coefficients}, \ref{thm:convergence} and \ref{thm:remainder} hold verbatim for the germ \eqref{eq:model} with $p<0$: the unique branch with $u\to0^+$ is $u=g(y)$ with $g\sim z$, $g^r$ has the expansion \eqref{eq:master} convergent for small $z$ (i.e.\ for $|y|$ large), and the remainder after the exponent cutoff $q$ in $w=1/|y|$ is $\OO(w^{(r+\sigma_*)/|p|}\M(w)^{k_*})$.
\end{theorem}
\begin{proof}
The derivative $f'(u)=apu^{p-1}(1+\oo(1))$ has a constant sign, so $f$ is monotone on $(0,\varepsilon)$ with $f\to\pm\infty$; the inverse exists on a neighbourhood of $\pm\infty$ and $u\sim z$. The proof of Theorem~\ref{thm:coefficients} is a computation at a fixed positive $q=v/a$ and never uses the sign of $p$. The parametric equation \eqref{eq:parametric} has $\partial_w\mathcal E(1,0)=p\ne0$ regardless of sign, and the specialisation $T_{ij}=z^{\delta_i}L^j\to0$ as $z\to0^+$ is the same. The remainder theorem is Lemma~\ref{lem:tail}.
\end{proof}

In the variable $y$ the blocks of $g^r$ read $(v/a)^{(r+\sigma)/p}C^{(r)}_\sigma\bigl(\tfrac1p\log(v/a)\bigr)$, with negative exponents of $y$ when $p<0$; the exponent of $w$ is $(r+\sigma)/|p|$, which is the quantity the cutoff refers to.

\subsection{Powers of the inverse and the original variable}

The observable power $r$ in \eqref{eq:master} gives directly the expansion of $u^r$. In terms of $x$:
\begin{equation}
\begin{aligned}
 (x-x_0)^r&=g(y)^r\ (x\to x_0^+),& (x_0-x)^r&=g(y)^r\ (x\to x_0^-),\\
 x^r&=g(y)^{-r}\ (x\to+\infty),& x^r&=(-1)^rg(y)^{-r}\ (x\to-\infty),
\end{aligned}
 \label{eq:power-back}
\end{equation}
The last equality in \eqref{eq:power-back}, involving $(-1)^r$, is used here only for integer $r$. Thus $x^r$ at infinity is obtained from \eqref{eq:master} with local-coordinate power $-r$. Both signed distances $x-x_0$ from above and $x_0-x$ from below are positive and admit arbitrary real powers. For $r=1$ and a finite endpoint, the original inverse is $x=x_0\pm g(y)$. A non-integer power of a negative original variable needs an additional branch convention and is outside the positive-real-power convention of this article.

\subsection{Examples}

\begin{example}[The Wright omega function]
$f(x)=x+\log x$ as $x\to+\infty$: $u=1/x$, $f=u^{-1}-\log u=u^{-1}(1-u\log u)$, so $a=1$, $p=-1$, one gap $\delta=1$ with $B=-L$, and target endpoint $+\infty$. The uniformizer is $z=1/y$, $L=-\log y$, and \eqref{eq:master} with $r=-1$ gives the original variable $x=1/g(y)$:
\begin{equation}
 x=y-\log y+\frac{\log y}{y}+\frac{\tfrac12\log^2y-\log y}{y^2}+\frac{\tfrac13\log^3y-\tfrac32\log^2y+\log y}{y^3}+\OO\Bigl(\frac{\log^4y}{y^4}\Bigr),
 \label{eq:omega}
\end{equation}
the classical expansion of the Wright omega function $\omega(y)=\W(e^y)$ \cite{cghjk,proveit}; the blocks are the Lambert polynomials $P_n(\log y)$ of \cite{proveit}.
\end{example}

\begin{example}[Reciprocal corrections]
$f(x)=x+1/x$ at $+\infty$: $u=1/x$, $f=u^{-1}(1+u^2)$, $p=-1$, gap $2$, constant coefficient. The inverse is the larger root of $x^2-yx+1=0$, $x=\bigl(y+\sqrt{y^2-4}\bigr)/2=y-y^{-1}-y^{-3}-2y^{-5}-\cdots$ (the coefficients are Catalan numbers), which \eqref{eq:constant-coefficients} reproduces.
\end{example}

\begin{example}[Half-integer gaps]
$f(x)=x+\sqrt x$ at $+\infty$: $f=u^{-1}(1+u^{1/2})$, gap $\frac12$. The inverse is $x=\bigl(\sqrt{4y+1}-1\bigr)^2/4=y-\sqrt y+\frac12-\frac1{8\sqrt y}+\frac1{128y^{3/2}}+\OO(y^{-5/2})$; the constant term $\frac12$ is the block of weight $2\cdot\frac12$, i.e.\ exponent $0$ in $w=1/y$, despite the infinite target endpoint.
\end{example}

\begin{example}[A finite endpoint from below]
$f(x)=x-x^2$ near $x_0=0$ with $x\to0^-$: $u=-x$, $f=-u-u^2$, $a=-1$, $p=1$, $y_0=0$, and $w=-y$ (the image is $y<0$). The inverse is $u=-y-y^2-2y^3-\cdots$, i.e.\ $x=y+y^2+2y^3+\OO(y^4)$, the branch of $x=(1-\sqrt{1-4y})/2$ continued to negative $y$.
\end{example}

\begin{example}[A pole]
$f(x)=1+1/x$ near $x_0=0^+$: $f=u^{-1}(1+u)$, $p=-1$, target endpoint $+\infty$, and $x=1/(y-1)=y^{-1}+y^{-2}+\cdots$ with remainder $\OO(y^{-3})$ after two terms.
\end{example}

\begin{example}[The second question, inverted]
For $F(x)=(1+x+x^{\sqrt2})^{\sqrt2}$ at $+\infty$ the forward germ is $u^{-2}(1+\sqrt2u^{\sqrt2-1}+\sqrt2u^{\sqrt2}+\cdots)$ from \eqref{eq:answer3}, with gaps generated by $\sqrt2-1$ and $\sqrt2$, $p=-2$, target endpoint $+\infty$, and $z=y^{-1/2}$. The inverse is
\begin{equation}
\begin{aligned}
 x={}&\sqrt y-\frac{1}{\sqrt2}\,y^{(2-\sqrt2)/2}+\Bigl(\frac34-\frac1{2\sqrt2}\Bigr)y^{(3-2\sqrt2)/2}\\
 &+\Bigl(1-\frac{5}{3\sqrt2}\Bigr)y^{(4-3\sqrt2)/2}-\frac{1}{\sqrt2}\,y^{(1-\sqrt2)/2}+\OO\bigl(y^{(5-4\sqrt2)/2}\bigr),
\end{aligned}
 \label{eq:second-inverse}
\end{equation}
the terms being ordered by decreasing exponent $-(r+\sigma)/2$ of $y$, with local-coordinate power $r=-1$. The forward remainder is transported by \eqref{eq:transport-power} with that value of $r$; the available forward order must be sufficient for the desired inverse cutoff. An independent derivation removes the outer power exactly: $1+x+x^{\sqrt2}=y^{1/\sqrt2}$, then applies the finite-germ coefficient formula.
\end{example}
